\documentclass[11pt]{article}
\usepackage[a4paper,margin=25mm]{geometry}
\usepackage{amsmath,amssymb,booktabs,hyperref}
\newcommand{\tht}{\vartheta}
\newcommand{\Q}{\mathbb Q}
\title{Exact boundary fixing for the quintic YY recursion\\A companion to version 0.2}
\author{}\date{16 October 2026}
\setlength{\emergencystretch}{2em}
\begin{document}\maketitle
\noindent
This supplement specifies the boundary layer of the runnable implementation.
The previous algebra conventions are retained. The new physical driver has
been executed through genus eight; genus 53 is its configured target, not a
claimed completed computation or performance result.

\section{Coordinate scale versus geometry}
Let $\mu\ne0$, $x=z/\mu=\psi^{-5}$, $\delta=1-x$, and
\[
 X=\frac{x}{x-1}=\frac1{1-\psi^5},\qquad
 \tht=\psi\partial_\psi=-5x\partial_x.
\]
The Picard--Fuchs operator is
\[
 L=\theta_x^4-x\prod_{j=1}^4(\theta_x+j/5),\qquad\theta_x=x\partial_x.
\]
Thus the MUM point is $x=0$, the conifold is $x=1$, and the orbifold is
$x=\infty$. In $z$ these are $0,\mu,\infty$. Leaving $\mu$ free changes the
coordinate scale, not the quintic exponents or topology. The standard quintic
coordinate is $z_{\rm can}=x/3125$.

For $f(x)=\sum c_nx^n$, the coefficient in the $z$ coordinate is
$c_n\mu^{-n}$. The implementation stores $c_n$ and the exponent $-n$
separately when this conversion is requested. It does not create a generic
$\Q(\mu)$ coefficient field. This removes unnecessary coordinate powers, but
not intrinsic rational or enumerative growth.

\section{Normalization and local holomorphic limits}
The original YY normalization is
\[
 C=\psi^3C_{\psi\psi\psi}=\frac{(2\pi i)^3}{125}(X-1),\qquad
 P_g=C^{g-1}F_g.
\]
In a local period frame $(\omega,t)$, the holomorphic generators are
\begin{align*}
 u&=\frac{\tht\omega}{\omega},&
 v_1&=\frac{\tht^2t}{\tht t}+2u,\\
 v_2&=\frac{\tht^2\omega}{\omega}-uv_1,&
 v_3&=\frac{\tht^3\omega}{\omega}+uv_2-5uX(v_1-2).
\end{align*}
The evaluated amplitude is
\[
 \mathcal F_g^{(\omega,t)}=
 \left[\frac{\omega^2(x-1)}{625}\right]^{g-1}P_g^{\rm hol}.
\]
Indeed the frame factor is $(2\pi i)^3\omega^2/(5^7C)$;
$5^7/5^3=5^4=625$, and $(X-1)^{-1}=x-1$.
All signs in the boundary equations below refer to this coupling convention.
A different string-coupling convention can change the written Bernoulli signs.

The anomaly primitive is fixed by $P_g^{\rm an}(0,0,0,X)=0$ and
\[
 P_g=P_g^{\rm an}+\sum_{j=0}^{3g-3}a_{g,j}X^j.
\]
Restrictions on a different Feynman representative cannot be imposed directly
on these raw $a_{g,j}$. We impose local conditions on the full evaluated amplitude.

\section{Local monodromy information}
The MUM Frobenius derivative basis without factorial division transforms by
$(T_0)_{ij}=\binom{i}{j}$ for $i\ge j$, with indices $0,\ldots,3$.
Consequently $(T_0-I)^4=0$ and $(T_0-I)^3\ne0$.
At the conifold, a logarithmic companion
$\nu\log\delta/(2\pi i)+\eta$ changes by $+\nu$ on a positive loop.
The nilpotent part has rank one and square zero.
The orbifold monodromy has eigenvalues $e^{2\pi i r/5}$,
$r=1,2,3,4$, and is represented by the companion matrix of
$\lambda^4+\lambda^3+\lambda^2+\lambda+1$.

These are three different local bases. We do not claim a shared integral
symplectic basis or a global connection matrix, and their matrices must not
be multiplied to check a global monodromy relation. The boundary solver uses
the local logarithmic structures together with Yukawa polarization and
enumerative topology. Monodromy by itself is insufficient.

\section{MUM series and constant maps}
The two required periods are $f$ and $f\log x+h$, where
\begin{align*}
 f&=\sum_{n\ge0}c_nx^n,&c_0&=1,&
 c_n&=c_{n-1}\frac{\prod_{j=1}^4(n-1+j/5)}{n^4},\\
 h&=\sum_{n\ge1}c_n\ell_nx^n,&
 \ell_n&=\sum_{k=1}^n\left(\sum_{j=1}^4\frac1{k-1+j/5}-\frac4k\right).
\end{align*}
The canonical instanton variable is
\[
 q=\frac{x}{3125}\exp(h/f)=\frac{z}{3125\mu}\exp(h/f).
\]
In the alternate variable $q_z=z\exp(h/f)$, a degree-$d$ canonical
coefficient would acquire $(3125\mu)^{-d}$. Canonical GV integrality therefore
must be imposed in $q$, not in $q_z$ for arbitrary $\mu$.

With $\chi=-200$ and the Bernoulli convention $B_1=-1/2$, the stable constant is
\[
 N_{g,0}=\frac{(-1)^g\chi\,|B_{2g}B_{2g-2}|}
 {2(2g)(2g-2)(2g-2)!}.
\]
It gives $N_{2,0}=-5/144$ and $N_{3,0}=5/36288$.
At $x=0$, all $v_i$ and $X$ vanish and the frame factor is $-1/625$;
therefore $a_{g,0}=(-625)^{g-1}N_{g,0}$.
For genus two this gives $a_{2,0}=3125/144$, including its positive sign.

\section{Conifold periods and the gap coefficient}
In $\delta$, $\theta_x=-(1-\delta)\partial_\delta$.
The analytic PF solution $f_c$ is specified by
$[\delta^0]f_c=1$, $[\delta]f_c=[\delta^2]f_c=0$.
These are a convenient rational A-period frame, not a declaration of an integral
period basis. The vanishing period is
\[
 \nu=\delta+\frac7{10}\delta^2+\frac{41}{75}\delta^3+\cdots,
 \qquad \tau=\nu/f_c.
\]
Here is how its second coefficient is fixed. Write
$L\delta^r=\sum_s A_s(r)\delta^{r+s}$.
Then
\[
 L(\delta^r\log\delta)=
 \sum_s\bigl(A_s(r)\log\delta+A_s'(r)\bigr)\delta^{r+s}.
\]
The coefficient of $\delta^{-1}$ in the nonlogarithmic part of
$L(\nu\log\delta)$ must vanish, since no analytic completion can cancel it.
Thus, for $\nu=\delta+b\delta^2+\cdots$,
$b=-A_{-2}'(1)/A_{-3}'(2)=7/10$.
The ordinary analytic recurrence fixes the remaining coefficients.
The implementation obtains these derivatives with exact dual numbers.

The genus-zero Yukawa in the standard $x$ gauge is
$C_{xxx}^{\rm can}=5/[x^3(1-x)]$.
Changing to $\delta$ gives a minus sign from three coordinate derivatives;
in the $f_c$ frame, near $\tau=0$, the residue is $-5/\tau$.
To obtain unit conifold residue, use $t_c=\sqrt5\,\tau$:
$C_{t_ct_ct_c}\sim-1/t_c$.
Consequently the gap in the rational frame is
\[
 \mathcal F_g^c(\tau)=
 \frac{5^{1-g}(-1)^{g-1}B_{2g}}{2g(2g-2)}
 \tau^{2-2g}+O(\tau^0).
\]
No irrational coefficients are required: $(\sqrt5)^{2-2g}=5^{1-g}$.
In particular the coefficients are $+1/1200$ at genus two and $+1/25200$
at genus three. The factor $5^{1-g}$ is a period normalization factor,
not a tunable extra condition.

The code subtracts this entire singular expression in $\delta$ and cancels
all negative-power coefficients. Their exponents range from $-2g+2$ to $-1$.
For example at genus two,
\[
 \frac1{1200}\tau^{-2}=
 \frac1{1200}\delta^{-2}-\frac7{6000}\delta^{-1}+O(1).
\]
Discarding the $\delta^{-1}$ term would impose an incorrect boundary condition.
Since $X=1-\delta^{-1}$ and the frame factor has valuation $g-1$,
there are $2g-2$ potentially negative powers.

\section{Orbifold regularity}
Let $y=\psi^5=1/x$ and
$\omega_r=\psi^r f_r(y)$, $r=1,2$, where
\[
 f_r(y)=\sum_{k\ge0}c_{r,k}y^k,\qquad c_{r,0}=1,\qquad
 \frac{c_{r,k}}{c_{r,k-1}}=
 \frac{(r+5k-5)^4}{\prod_{j=1}^4(r+5k-j)}.
\]
These follow by substituting into the original psi PF operator.
Take $\omega=\omega_1$, $t=\omega_2/\omega_1$.
Writing $D=5y\partial_y$ gives
\begin{align*}
 u&=1+D\log f_1,&
 T&=1+D\log(f_2/f_1),\\
 v_1&=T+D\log T+2u,&
 B_2&=Du+u^2,&B_3&=DB_2+uB_2.
\end{align*}
The explicit $1$ in $u$ and $T$ comes from the stripped psi powers.
Use $v_2=B_2-uv_1$ and $v_3=B_3+uv_2-5uX(v_1-2)$.
Since
\[
 \omega_1^2(x-1)/625=\psi^{-3}f_1(y)^2(1-y)/625,
\]
regularity means that the coefficients of
\[
 S_g(y)=\left[f_1(y)^2(1-y)/625\right]^{g-1}P_g^{\rm hol}(y)
\]
vanish for $0\le k<\lceil(3g-3)/5\rceil$.
This removes negative psi powers in $\psi^{-3g+3}S_g(\psi^5)$.

\section{Enumerative input and genus 53}
For positive degree, use
\[
 \sum_{g\ge0,d\ge1}N_{g,d}\lambda^{2g-2}q^d=
 \sum_{h\ge0,d\ge1,r\ge1}\frac{n_h^d}{r}
 \bigl(2\sin(r\lambda/2)\bigr)^{2h-2}q^{rd},
\]
where the left genus sum also includes genus zero.
To remove any possible indexing ambiguity, its coefficient form is
\[
 N_{g,d}=\sum_{r\mid d}r^{2g-3}\sum_{h=0}^g
 n_h^{d/r}\,[w^{g-h}]
 \left(\sum_{j\ge0}\frac{(-1)^j w^j}{2^{2j}(2j+1)!}\right)^{2h-2}.
\]
The sine-series constant is one. Thus the $r=1,h=g$ coefficient is exactly
$n_g^d$; proper covers and lower genera supply the other terms.

The proven Castelnuovo input is $n_g^d=0$ when
$10g>d^2+5d+10$.
Through genus 53, all same-genus invariants needed for
$d\le D_g=\lfloor(2g-2)/5\rfloor$ vanish except
\[
 n_{51}^{20}=175.
\]
For degree 20 the extremal genus is 51. Liu--Ruan's extremal formula, with
$m=4$, gives $(-1)^{\binom73-\binom23+3}5(\binom73-\binom23)
=(-1)^{38}\cdot175=175$.
In contrast $n_{52}^{20}=n_{53}^{20}=0$ by the strict vanishing bound.
Unknown invariants outside the supported range are never silently set to zero.

The genus-zero Yukawa and genus-one derivative seed initialize the GV table;
lower-genus computed amplitudes provide subsequent entries. To avoid explicit
series reversion, for $q=x\exp(m(x))/3125$ use Lagrange inversion:
\[
 [q^d]A(x(q))=\frac{3125^d}{d}
 [x^{d-1}]A'(x)\exp(-dm(x)),\qquad d\ge1.
\]
The MUM equations can instead compare power series directly in $x$.

There are $3g-2$ raw pure-$X$ coefficients and
\[
 \underbrace{\left\lceil\frac{3g-3}5\right\rceil}_{\rm orbifold}
 +\underbrace{(2g-2)}_{\rm conifold}
 +\underbrace{\left(1+\left\lfloor\frac{2g-2}5\right\rfloor\right)}_{\rm MUM}
 =3g-2
\]
conditions. Actual exact rank is checked as well as the count.
At genus 53 these numbers are $32+104+21=157$ and $D_{53}=20$.
These are boundary linear-system sizes, not sizes of the anomaly polynomials.

\section{Certification, checkpoints, and practical limits}
A Laurent jet stores its absolute precision $O(t^N)$.
For valuations $a,b$, multiplication has precision
$\min(N_f+b,N_g+a)$; inversion has precision $N_f-2a$.
Uncertified coefficient access raises. A conservative working precision
$6g+16$ is used and low-genus answers are tested again at larger precision.
Nested Horner evaluation reduces repeated jet products.
Polynomial and jet multiplication collect common denominators and multiply
integers internally, with rational reduction at output coefficients.

The physical driver checks the full BCOV residual, exact boundary rank,
GV integrality, and available Castelnuovo values before committing a genus.
On resume it recomputes the physical boundary vector as well as the anomaly;
a formally valid zero-ambiguity checkpoint is rejected.
Resource-limit errors retain the last completed physical checkpoint.
The standard-library expanded-polynomial backend is not claimed to be practical
at genus 53. Leaving $\mu$ free removes coordinate artifacts but does not
replace a high-performance polynomial kernel or certified modular reconstruction.

\begin{thebibliography}{9}
\bibitem{YY} S. Yamaguchi and S.-T. Yau,
\emph{Topological String Partition Functions as Polynomials},
\href{https://arxiv.org/abs/hep-th/0406078}{hep-th/0406078}.
\bibitem{HKQ} M.-x. Huang, A. Klemm, and S. Quackenbush,
\emph{Topological String Theory on Compact Calabi--Yau: Modularity and Boundary Conditions},
\href{https://arxiv.org/abs/hep-th/0612125}{hep-th/0612125}.
\bibitem{LR} Z. Liu and Y. Ruan,
\emph{Castelnuovo bound and higher genus Gromov--Witten invariants of quintic 3-folds},
\href{https://arxiv.org/abs/2210.13411}{arXiv:2210.13411v2},
Theorems 1.1, 1.3 and Appendix A.
\end{thebibliography}
\end{document}
